\section{Introduction}
\subsection{Goals.}
I study the $\pi^0$ electroproduction of the $\Delta(1232)$ resonance in the momentum transfer range
$Q^2:\, 2\,\rightarrow\,5.5\,\,GeV^2$ produced in Hall B at  Jefferson Laboratory
by an electron beam having an energy of $5.7542\,GeV$ scattering off a liquid hydrogen target.

The four-fold cross section will be extracted as a function of $W\, ,Q^2\, ,\phi^*$ and $cos\theta^*$.
The $P\rightarrow \Delta$ magnetic form factor $G_M^*$ and the ratios of electric and Coulomb
electromagnetic quadrupoles to the magnetic dipole ($R_{EM}$ and $R_{SM}$) will therefore be extracted.

\subsection{Motivations.}
The constituent quark model (CQM) describes rather well the static ($Q^2=0$) property of the nucleons.
It couples $(u,d,s)$ flavor symmetry $SU(3)$ with spin symmetry $SU(2)$ to give a $SU(6)$ framework
which yelds the following multiplet structure:
$$6\otimes 6\otimes 6 = 56_S \oplus 70_{MS} \oplus 70_{MA} \oplus 20_{A} $$ 
where the subscript denotes the symmetry under interchange of two of the three quarks and the asymmetry
of the total wave function is granted by color $SU(3)_C$. 
In particular, the
proton and the $\Delta(1232)$ resonance belong to the same multiplet with $L=0$ and parity $P=(-1)^L=+1$
which is $(56_S, 0^+)$, where the second number is $L^P$. 
The transition $P\rightarrow\Delta^+$ does not change the angular momentum 
 but it flips the spin of one of the quarks:
$$
u\uparrow d\uparrow u\downarrow\,\,\longrightarrow\,\, u\uparrow d\uparrow u\uparrow
$$
such that both the isospin and the spin increase by a unit, i.e. 
$$I:\,\frac{1}{2}\rightarrow\frac{3}{2}  \qquad  J:\,\frac{1}{2}\rightarrow\frac{3}{2}    .$$
Hence this is a pure magnetic transition with $\Delta J=+1$, denoted $M1$.

In its simplest form, the CQM prescribes a harmonic oscillator potential in which the quarks are moving
non relativistically. This accounts for several baryon properties, for example the $(56, 0^+),(70, 1^-)$
energy splitting. However the $P$ and the $\Delta$ would have the same mass.
Various improvements has been introduced over the past few decades by introducing two body interactions \cite{bib:twobody},
relativizing the basic Hamiltonian \cite{bib:relatiH}, adding a quark-quark hyperfine interaction \cite{bib:hyperfine},
 and including gluonic degrees of freedom \cite{bib:gluonic}.
The hyperfine interaction accounts for the $P\, -\,\Delta$ energy splitting and introduces a contribution of $L=2$ to the $s$ wave
functions of $P$ and $\Delta$, which is observed experimentally. This introduces small quadrupoles ($L=2$) 
moments\footnote{The number $2$ refers to the angular content of the electromagnetic field of the
real or virtual photon interacting with the nucleon (leptonic vertex). One could also refer to the angular 
contents of the resonance (hadronic vertex) by means of both electromagnetic and strong force, in 
which case $M1$ is written as $M_{1+}$, $E2$ is written as $E_{1+}$ and $C2$ is written as $S_{1+}$.}
$E2$ (electric) and $C2$ (Coulomb).
However, as $Q^2$ increases the CQM loses predictive power and other models are required. 

Quantum Chromo Dynamics can be used in the regime where the coupling constant $\alpha$ is small using the perturbation
approach (pQCD). In this case helicity conservation requires $E2 = M1$  and $C2 = 0$, an asymptotic behavior
not observed for $Q^2$ up to $4\, GeV^2$ which is at the moment the maximum measured value of $Q^2$.

Numerous models exist in the so called CQM $\rightarrow$ QCD transition regime that predict the magnetic moments
of the $N \,- \,\Delta$ transition. Among the others:

\begin{itemize}
\item {\it a diquark model} \cite{bib:diquark}. In this model the baryon consists of a quark and a couple of correlated
      pair of quarks called ``diquark``. It allows for helicity violation at low $Q^2$ and requires its conservation
      at high $Q^2$.  

\item {\it an effective lagrangian} \cite{bib:efflag1}, \cite{bib:efflag2}. In this approach $P\, ,\Delta\, ,\pi\, ,\gamma$ (and not
      the quark and gluons) are the degree of freedom and 
      the used effective lagrangian $L=L^1_{\pi N\Delta}+L^2_{\gamma N\Delta}$ incorporates chiral symmetry.

\item {\it a unitary isobar model} \cite{bib:maid2000}. This model prescribes a new operator for pion production which includes
      Born term, vector meson and nucleon resonances. The model is developed in Mainz and it is known as {\bf maid2000}.

\item {\it the Sato Lee model} \cite{bib:pion2}, \cite{bib:sato1}, \cite{bib:sato2}. A dynamical model 
      with intermidiate rescattering terms, developed 
      specifically for the $\Delta(1232)$ resonance. 
      It includes five terms in the hamiltonian: a free hamiltonian $H_0$, a non-resonant $\pi N$ potential $v_{\pi N}$, 
      a non resonant $\gamma N\rightarrow \pi N$ transition $v_{\gamma N}$ and the two excitation modes $\Gamma_{\gamma N \leftrightarrow\Delta}$
      and $\Gamma_{\pi N \leftrightarrow\Delta}$.

\item {\it a dynamical model} \cite{bib:dmt2001}. This model is derived from maid2000, but in this
      case the background amplitudes do not ``dress'' the $\gamma P\Delta$ vertex
      and the helicity amplitudes correspond to the ``bare'' vertex. 
      The model is developed in Mainz and it is known as {\bf DMT model}, or {\bf maid2001}.

\item {\it a chiral quark-soliton model} \cite{bib:soli}. In this model the lagrangian 
      $${\cal L }= \overline{\Psi}(\imath\gamma^\mu\partial_\mu - m - MU^{\gamma_5})\Psi$$ contains a spontaneous chiral symmetry breaking
      dynamical quark mass $M$.



\end{itemize}
\newpage




